
\begin{tabularx}{\textwidth}{@{}lX@{}}
\toprule
Dato & Valor \\
\midrule
CADi & IA agéntica para la creación de contenido académico \\
Escuela & Escuela de Ingeniería y Ciencias (EIC) \\
Clave sugerida & \clave \\
Duración & \textbf{40 horas totales}: 24\,h en aula (8 $\times$ 3\,h) + 16\,h fuera de aula \\
Modalidad & Virtual (Teams) / híbrida / presencial según edición \\
Audiencia & Profesorado de profesional y posgrado de la EIC \\
Instructor & Juliho Castillo \\
Idioma & Español (México) \\
Colaboración & CEDDIE / oferta de CADi de la Escuela \textit{(verificar por edición)} \\
\bottomrule
\end{tabularx}

\section{Descripción}

Este CADi forma al profesorado de la EIC en el diseño y operación de \textbf{flujos de trabajo con IA agéntica} para producir contenido académico de calidad típico de ingeniería y ciencias. Se trabaja con roles de agente, herramientas, orquestación y supervisión humana, integrando criterios de integridad académica (verificación de ecuaciones, citas y datos) y plantillas reutilizables.

El formato combina \textbf{8 sesiones de 3 horas en aula} (24\,h) con \textbf{16 horas de trabajo independiente / asíncrono}.

\section{Objetivos generales}

\begin{enumerate}
\item Comprender el paradigma agéntico aplicado a materiales académicos STEM.
\item Diseñar y ejecutar pipelines asistidos (brief $\rightarrow$ orquestación $\rightarrow$ generación $\rightarrow$ revisión $\rightarrow$ calidad).
\item Transferir plantillas y prácticas al contexto docente e investigativo en la EIC.
\item Documentar de forma transparente el uso de IA (agent log) y la autoría humana.
\item Completar un \textbf{microartefacto EIC} con evidencia de proceso a lo largo de 40 horas.
\end{enumerate}

\section{Resultados de aprendizaje}

Ver tabla RA1--RA8 en la propuesta. En resumen: fundamentos e integridad; brief, prompt y orquestación; arquitectura/outline de libro; secciones y consistencia; narrativa/storyboard; deck y notas; estructura/borrador de artículo; revisión agéntica y cierre.

\section{Estructura por sesión (aula --- 24\,h)}

\subsection{Sesión 1 --- Fundamentos agénticos e integridad académica (EIC) (3\,h)}
\textbf{Objetivos:} Distinguir agentes de chat; enmarcar integridad en contextos EIC; identificar riesgos (ecuaciones, citas, datos).\\
\textbf{Temas:} Definiciones; herramientas y memoria; HITL; riesgos; política de transparencia (\ceddie).\\
\textbf{Producto (E1):} Mapa conceptual agentes vs.\ chat + nota de integridad (5--8 líneas) + tema EIC tentativo.

\subsection{Sesión 2 --- Briefs, prompts de sistema y orquestación (3\,h)}
\textbf{Objetivos:} Diseñar brief y prompt reutilizables; orquestación multi-rol con puntos HITL.\\
\textbf{Producto (E2):} Brief + prompt de sistema v0 + esquema de orquestación (2--4 roles).

\subsection{Sesión 3 --- Libros/ebooks I --- Arquitectura y outline (3\,h)}
\textbf{Producto (E3):} Arquitectura de obra + outline documentado (6--10 nodos) con RA por capítulo/sección clave.

\subsection{Sesión 4 --- Libros/ebooks II --- Secciones, ejercicios, consistencia (3\,h)}
\textbf{Producto (E4):} 1--2 secciones borrador + artefacto didáctico + reporte breve de consistencia.

\subsection{Sesión 5 --- Presentaciones I --- Objetivos, narrativa, storyboard (3\,h)}
\textbf{Producto (E5):} Storyboard narrativo + justificación de mensajes clave.

\subsection{Sesión 6 --- Presentaciones II --- Deck, notas, accesibilidad (3\,h)}
\textbf{Producto (E6):} Deck MV (8--12 slides) + notas de orador + mini-checklist de accesibilidad.

\subsection{Sesión 7 --- Artículos I --- Estructura, borrador, citas (3\,h)}
\textbf{Producto (E7):} Estructura + borrador de sección/micro-artículo con lista de citas/datos a verificar.

\subsection{Sesión 8 --- Artículos II --- Revisión, cierre, showcase (3\,h)}
\textbf{Producto de acreditación (E8):} Microartefacto EIC + agent log + checklist de calidad.\\
Opciones: capítulo corto STEM, deck de unidad con notación técnica, sección tipo IMRyD, o guía de laboratorio breve.

\section{Trabajo fuera de aula (16\,h)}

\begin{tabularx}{\textwidth}{@{}Xcccc@{}}
\toprule
Bloque & Horas & Sesiones & Actividades & Evidencia \\
\midrule
A. Lecturas y demos & 4\,h & 1--2 & Lecturas; demos del stack & Bitácora \\
B. Práctica stack/plantillas & 4\,h & 2--6 & Brief, prompt, checklist & Plantillas iteradas \\
C. Avance microartefacto & 6\,h & 3--8 & Outline$\rightarrow$secciones / deck / borrador & Versiones fechadas \\
D. Agent log + checklist & 2\,h & 7--8 & Log, checklist, reflexión & Paquete de cierre \\
\textbf{Total} & \textbf{16\,h} & --- & --- & --- \\
\bottomrule
\end{tabularx}

\begin{quote}
El facilitador puede reasignar $\pm$0{,}5\,h entre bloques A--C según cohorte, siempre que el total fuera de aula sea \textbf{16\,h}.
\end{quote}

\section{Metodología}

Cada sesión en aula combina: (1) marco conceptual breve; (2) demostración en vivo; (3) práctica guiada con plantillas; (4) puesta en común; (5) encargo de TI.

\section{Evaluación}

\begin{tabularx}{\textwidth}{@{}lXl@{}}
\toprule
Evidencia & Descripción & Momento \\
\midrule
E1 & Mapa agentes vs.\ chat + nota de integridad + tema EIC & 1 \\
E2 & Brief + prompt + orquestación & 2 \\
E3 & Arquitectura + outline & 3 \\
E4 & Secciones + ejercicio/figura + consistencia & 4 \\
E5 & Storyboard narrativo & 5 \\
E6 & Deck MV + notas + accesibilidad & 6 \\
E7 & Estructura + borrador de artículo/sección & 7 \\
E8 & Microartefacto + agent log + checklist & 8 \\
TI & Evidencia de las 16\,h independientes & Continuo \\
\bottomrule
\end{tabularx}

\begin{quote}
\ceddie\ el umbral oficial de acreditación.
\end{quote}

\section{Integridad académica y uso responsable de IA}

\begin{enumerate}
\item La persona docente es \textbf{autora responsable} del contenido publicado o impartido.
\item Toda asistencia sustancial de IA debe poder \textbf{documentarse} (agent log).
\item Afirmaciones factuales, citas, ecuaciones y datos requieren \textbf{verificación humana}.
\item No se deben ingresar datos personales sensibles ni material confidencial no autorizado.
\item Se declara el uso de IA según las normas del medio de publicación o del curso.
\end{enumerate}

\begin{quote}
\ceddie\ los textos oficiales a citar en convocatoria y en el aula.
\end{quote}

\section{Contacto}

Instructor: Juliho Castillo --- \textit{(correo institucional a completar)}.\\
EIC · Coordinación académica / CEDDIE: \textit{(contacto de la edición)}.
